% ============================================================
% DOCUMENTCLASS
% ============================================================
\documentclass[12pt,oneside]{book}

% ============================================================
% METADATOS
% ============================================================
\newcommand{\universidad}{Universidad de Buenos Aires}
\newcommand{\facultad}{Facultad de Derecho}
\newcommand{\materia}{Derecho de la Integración}
\newcommand{\titulo}{Parte 2: Derecho laboral en el Mercosur y la Unión Europea}
\newcommand{\subtitulo}{Consultas y repaso para el parcial}
\newcommand{\autor}{Isaac Edgar Camacho Ocampo}
\newcommand{\anio}{2026}

% ============================================================
% PAQUETES
% ============================================================
\usepackage[utf8]{inputenc}
\usepackage[T1]{fontenc}
% En esta instalación no está disponible el soporte de babel para español; se omite para
% que el documento compile correctamente y genere el índice en el PDF.

\usepackage[
    top=2.5cm,
    bottom=2.5cm,
    left=3cm,
    right=2.5cm,
    headheight=15pt
]{geometry}

\usepackage{amsmath}
\usepackage{amssymb}
\usepackage{mathtools}

\usepackage{graphicx}
\usepackage{float}
\usepackage{caption}
\usepackage{subcaption}

\usepackage{booktabs}
\usepackage{array}
\usepackage{longtable}
\usepackage{tabularx}

\usepackage{enumitem}

\usepackage{xcolor}
\usepackage[most]{tcolorbox}

\usepackage{fancyhdr}
\usepackage{ragged2e}

\usepackage[
    colorlinks=true,
    linkcolor=blue!50!black,
    citecolor=green!40!black,
    urlcolor=blue!60!black,
    pdftitle={Apuntes de Derecho de la Integración - Parte 2},
    pdfauthor={Isaac Edgar Camacho Ocampo},
    pdfsubject={Apuntes universitarios},
    bookmarks=true
]{hyperref}

% ============================================================
% CONFIGURACIÓN
% ============================================================
\graphicspath{
    {./img/}
    {./graficos/}
    {./figuras/}
}

\setcounter{tocdepth}{2}

\setlength{\parindent}{0pt}
\setlength{\parskip}{0.5em}

\setlist[itemize]{
    topsep=0.3em,
    itemsep=0.2em
}

\setlist[enumerate]{
    topsep=0.3em,
    itemsep=0.2em
}

\clubpenalty=10000
\widowpenalty=10000

\newcolumntype{L}[1]{>{\RaggedRight\arraybackslash}p{#1}}
\newcolumntype{Y}{>{\RaggedRight\arraybackslash}X}

% ============================================================
% COMANDOS
% ============================================================
\newcommand{\abs}[1]{\left|#1\right|}
\newcommand{\norm}[1]{\left\lVert#1\right\rVert}
\newcommand{\vect}[1]{\mathbf{#1}}
\newcommand{\deriv}[2]{\frac{d#1}{d#2}}
\newcommand{\pderiv}[2]{\frac{\partial#1}{\partial#2}}

% ============================================================
% ENTORNOS / CAJAS
% ============================================================
\newtcolorbox{definition}[1]{
    enhanced,
    breakable,
    title={Definición: #1},
    fonttitle=\bfseries,
    colback=blue!3,
    colframe=blue!50!black,
    boxrule=0.8pt,
    arc=2mm
}

\newtcolorbox{important}{
    enhanced,
    breakable,
    title={Importante},
    fonttitle=\bfseries,
    colback=yellow!8,
    colframe=orange!70!black,
    boxrule=0.8pt,
    arc=2mm
}

\newtcolorbox{example}[1]{
    enhanced,
    breakable,
    title={Ejemplo: #1},
    fonttitle=\bfseries,
    colback=green!4,
    colframe=green!40!black,
    boxrule=0.8pt,
    arc=2mm
}

\newtcolorbox{formula}[1]{
    enhanced,
    breakable,
    title={Fórmula / Regla: #1},
    fonttitle=\bfseries,
    colback=purple!4,
    colframe=purple!50!black,
    boxrule=0.8pt,
    arc=2mm
}

\newtcolorbox{exercise}[1]{
    enhanced,
    breakable,
    title={Ejercicio: #1},
    fonttitle=\bfseries,
    colback=gray!5,
    colframe=gray!60!black,
    boxrule=0.8pt,
    arc=2mm
}

\newtcolorbox{note}{
    enhanced,
    breakable,
    title={Nota},
    fonttitle=\bfseries,
    colback=white,
    colframe=gray!50,
    boxrule=0.6pt,
    arc=2mm
}

% ============================================================
% CONFIGURACIÓN FINAL
% ============================================================
\pagestyle{fancy}
\fancyhf{}

\fancyhead[L]{\nouppercase{\leftmark}}
\fancyhead[R]{\materia}

\fancyfoot[C]{\thepage}

\renewcommand{\headrulewidth}{0.4pt}
\renewcommand{\footrulewidth}{0pt}

\fancypagestyle{plain}{
    \fancyhf{}
    \fancyfoot[C]{\thepage}
    \renewcommand{\headrulewidth}{0pt}
}

% ============================================================
% DOCUMENT
% ============================================================
\begin{document}

% ------------------------------------------------------------
% PORTADA
% ------------------------------------------------------------
\begin{titlepage}

\thispagestyle{empty}

\begin{center}

\vspace*{1cm}

{\large \textsc{\universidad}}

\vspace{0.2cm}

{\large \textsc{\facultad}}

\vspace{3cm}

{\Huge \textbf{\materia}}

\vspace{1cm}

{\LARGE \titulo}

\vspace{0.4cm}

{\large \subtitulo}

\vspace{0.8cm}

\textit{Estudiar con constancia convierte cada concepto en conocimiento.}

\vspace{1.2cm}

\rule{0.70\textwidth}{0.5pt}

\vspace{0.5cm}

{\large Apuntes de estudio}

\vspace{0.5cm}

\rule{0.70\textwidth}{0.5pt}

\vfill

{\large \autor}

\vspace{0.3cm}

{\large \anio}

\end{center}

\vfill

\begin{flushleft}
\textit{Este es un modesto aporte a los estudiantes de ``\materia'' no pretende sustituir el material oficial de la materia.}
\end{flushleft}

\end{titlepage}

% ------------------------------------------------------------
% ÍNDICE
% ------------------------------------------------------------
\tableofcontents

% ============================================================
% CAPÍTULO 1
% ============================================================
\chapter[Valores, trabajo e igualdad]{Introducción: valores compartidos, concepto de trabajo e igualdad de trato}

% Se omitieron digresiones políticas, religiosas y anecdóticas de la charla inicial sin valor conceptual para la materia.

\section{Punto de partida: los valores comunes}

El desarrollo parte de la reflexión sobre la visita del Papa a Latinoamérica, utilizada como puente hacia el tema central: el \textbf{derecho al trabajo y los derechos sociales dentro de los espacios de integración}.

Se plantea la pregunta de qué se fue a buscar en esa visita y por qué concurrió tanta gente, incluso personas no practicantes. Se sostiene que, más allá de lo religioso, lo buscado fue una palabra de aliento y de esperanza (la idea de que, a pesar de las caídas, es posible levantarse y seguir adelante) que resulta común a los argentinos y que constituye un punto de unión por encima de las diferencias políticas.

Entre las lecturas expresadas sobre el sentido de la visita al continente se mencionaron:
\begin{itemize}
    \item la búsqueda de unión y de no separarse tanto;
    \item la visibilización de la acción social que las iglesias realizan en los barrios (comedores, Cáritas, asistencia a personas en situación de calle, a enfermos y a ancianos);
    \item la atención hacia personas con discapacidad, jubilados y niños;
    \item la fuerte presencia del catolicismo en Latinoamérica (se mencionó que aproximadamente el 80\,\% de la población es católica).
\end{itemize}

Los valores que se señalan como compartidos y que ocupan la agenda social son la protección de los jóvenes, la protección de los ancianos, la lucha contra las drogas y la protección de las personas con discapacidad. Si tantas personas coinciden en ello, se concluye que hay más en común que diferencias y que «no tenemos que pelearnos tanto»: esa base común permite trabajar en conjunto.

La conexión con la materia es la siguiente: el tema es el \textbf{derecho laboral dentro del ámbito regional} (el Mercosur, comparado con la Unión Europea), es decir, los derechos sociales en un ámbito de integración: igualdad de trato, derecho al trabajo, protección social y descanso.

\section{El concepto de trabajo}

Se observa que suele entenderse como trabajo solamente el que es remunerado, que es el que la legislación considera. De esa noción proviene la visión de que el trabajo de ama de casa «no es trabajo».

Frente a ello se propone incluir dentro de la idea de \textbf{labor} el \textbf{trabajo social}: quien presta un servicio a otro realiza trabajo social. En ese sentido, las amas de casa y las organizaciones brindan servicio a la educación y al bienestar de otras personas (la familia, los hijos) y forman la mente y la conducta de mujeres y varones; se trata de un trabajo digno y necesario.

Se advierte además un \textbf{doble discurso}: la misma tarea es valorada cuando se realiza en una institución pública y desvalorizada cuando se realiza en el hogar. También se señala que el discurso según el cual la mujer fue relegada a las tareas del hogar y no trabaja si no está inserta en las oficinas desvaloriza tareas vinculadas con la organización social.

Se observa que las posturas de izquierda y de derecha se atribuyen mutuamente la responsabilidad: una acusa a la otra de medirlo todo en valor económico, y la otra le atribuye la discriminación y la segregación de la mujer. Como explicación histórica de la distribución de tareas (trabajo fuera o dentro del hogar), se mencionan razones sociales: la mujer debía amamantar y el hombre tenía mayor fuerza física, de modo que el trabajo pesado recaía en el hombre y el trabajo de la casa, más liviano, en la mujer.

En las intervenciones se planteó que lo que se valora es el trabajo vinculado con lo masculino, y que para que el trabajo propio sea valorado no debería ser necesario hacer lo que hace el otro: debería valorarse el cuidado de otra persona como un valor social.

\section{Igualdad de trato, cupos y acceso a la capacitación}

\begin{example}{la terna de jueces y el cupo femenino (experiencia en el Consejo de la Magistratura)}
Se aprobó una fundamentación según la cual en toda terna para la elección de jueces debía haber siempre una mujer, como garantía del cupo femenino. El argumento era que la mujer tiene menos oportunidades de profesionalizarse y perfeccionarse (por las tareas de cocina y crianza de los hijos), por lo que haría menos cursos de posgrado y obtendría siempre menor puntaje que el hombre en un examen.

Frente a ello se expuso que se trata de una \textbf{premisa falsa}: una mujer libre elige casarse y quedarse en su casa a criar a su hijo, y no tiene por qué pagar las consecuencias el justiciable al que se le nombre un juez que no se profesionalizó. Con ese criterio, el candidato que salió tercero en el concurso debería ser desplazado para que ingrese la mujer que salió décima. Se agrega que el justiciable debe aspirar al mejor juez posible y que el planteo del cupo implica afirmar que la mujer no tiene capacidad para competir en igualdad de condiciones.
\end{example}

La \textbf{Declaración Sociolaboral del Mercosur} garantiza la igualdad de trato en el acceso al trabajo entre hombres y mujeres, sin discriminación por razón de religión, credo, raza o sexo. Sobre esa base se discute si garantizar un cupo equivale a tratar a todos por igual.

\begin{note}
\textbf{Para reflexionar:} se plantea la pregunta de si garantizar el cupo es tratar igual. La respuesta que se expone es negativa.
\end{note}

Las posiciones expuestas son las siguientes:
\begin{itemize}
    \item \textbf{Cupo como herramienta mínima.} En una intervención se sostuvo que el cupo es una herramienta mínima, que puede usarse, pero no con la justificación de una entidad de género; se la considera exageradamente usada, y el problema se origina en las oportunidades desde el inicio de la vida.
    \item \textbf{La igualdad de trato pasa por el acceso a la capacitación.} Si se propicia y facilita el acceso a capacitación (por ejemplo, cursos laborales a distancia que hagan razonar y no cursos que otorguen un título sin contenido), luego se puede competir en igualdad de condiciones por un cargo jerárquico, judicial o gerencial.
    \item \textbf{Extensión del criterio del cupo.} Con el criterio del cupo habría que establecer también cupos para personas afrodescendientes, con discapacidad, etc. Se agrega que las mujeres no son una minoría, sino la mitad del mundo.
\end{itemize}

\begin{example}{la postura de la profesora Hebe Leonardi de Herbón sobre el cupo femenino}
Cuando la docente era estudiante, la profesora Hebe Leonardi de Herbón (feminista) sostenía que el cupo femenino no era conveniente porque parte de la base de que la mujer no puede competir en igualdad de condiciones con los hombres, lo que sería reconocer que se depende de ellos para estar iguales. Aun así, admitía que en una primera instancia de lucha era necesario llegar a ese primer escalón para luego estar en posición de luchar. Pasados 40 años se sigue luchando por el cupo femenino, lo que lleva a concluir que ese camino no está resultando: se critica que el cupo se cubra con personas no capacitadas para la función que ejercen. La alternativa propuesta es el \textbf{acceso a la capacitación}.
\end{example}

% ------------------------------------------------------------
% CORNELL CAPÍTULO 1
% ------------------------------------------------------------
\section{Cuadro Cornell del capítulo}

\begin{longtable}{L{0.27\textwidth}L{\dimexpr0.73\textwidth-4\tabcolsep\relax}}
\toprule
\textbf{Preguntas / palabras clave} & \textbf{Notas / respuestas} \\
\midrule
\endfirsthead
\toprule
\textbf{Preguntas / palabras clave} & \textbf{Notas / respuestas} \\
\midrule
\endhead

\textbf{¿Por qué se parte de la visita del Papa y qué tiene que ver con la integración?} &
Es el disparador: la multitud mostró valores compartidos (protección de jóvenes y ancianos, lucha contra la droga, asistencia a personas con discapacidad). Si hay tanto en común, «no tenemos que pelearnos tanto»; de ahí se pasa al derecho al trabajo y a los derechos sociales en los espacios de integración. \\[0.6em]

\textbf{¿Qué se entiende por «trabajo» más allá del trabajo remunerado?} &
Incluye el trabajo social y las tareas del hogar. La legislación mira el trabajo remunerado, pero existe un «doble discurso» que valora la tarea realizada en una institución y desvaloriza la misma tarea realizada en la casa. \\[0.6em]

\textbf{¿Qué postura se expuso sobre los cupos y la igualdad de trato?} &
El cupo es una herramienta mínima y exageradamente usada. La igualdad de trato pasa por el acceso a la capacitación y por la igualdad de oportunidades desde el inicio de la vida. \\

\midrule
\multicolumn{2}{L{\dimexpr\textwidth-2\tabcolsep\relax}}{\textbf{Resumen:} la clase parte de la visita del Papa para mostrar que, más allá de las diferencias políticas, existen valores sociales compartidos que sostienen cualquier proceso de integración, y desemboca en el derecho al trabajo. Se amplía el concepto de trabajo más allá del remunerado y se sostiene que la igualdad de trato depende del acceso a la capacitación más que de los cupos.} \\
\bottomrule
\end{longtable}

% ============================================================
% CAPÍTULO 2
% ============================================================
\chapter[Derecho laboral: Mercosur]{Derecho laboral en el Mercosur}

\section{Marco normativo: la Declaración Sociolaboral y su técnica de redacción}

En todo ámbito de integración, y fundamentalmente en los de tinte económico y social, la libertad de trabajo y las cuestiones socioeconómicas hacen al desarrollo del espacio de integración. Por ello deberían existir normas específicas del ámbito de integración que regulen estas cuestiones de manera uniforme o, al menos, armónica.

La norma fundamental del Mercosur en la materia es la \textbf{Declaración Sociolaboral del Mercosur}, elaborada por la comisión de asuntos sociolaborales que trabajaba dentro del Grupo Mercado Común. Se recuerda que el Grupo Mercado Común trabaja por comités y subgrupos de trabajo, y que existe un foro económico y social integrado por delegados de consumo, de trabajo y empresariales. La Declaración es del año 2015 y fue firmada por los presidentes.

Su redacción adopta una técnica determinada: «Los Estados se comprometen a\ldots». La tarea asignada a la comisión consiste en elaborar informes y estadísticas y hacer propuestas al Grupo Mercado Común sobre normativas o cambios, y toda propuesta debe darse por consenso. La comisión debe reunirse al menos dos veces al año para presentar los informes sobre lo trabajado.
% TODO: contenido incompleto o ambiguo en las notas originales (se indica que el artículo prevería una reunión anual, pero que actualmente serían dos al año).
No hay una regulación concreta: la Declaración se mantiene en aquello en lo que todos están de acuerdo, como la \textbf{no discriminación} por razones de raza, religión o sexo. Se critica que en 2015 hiciera falta una declaración firmada por los presidentes para establecer como compromiso de los Estados algo tan básico.

\begin{important}
La Declaración es una \textbf{declaración}, no una convención. Se la caracteriza como una expresión de buenos deseos: declara que se garantizan derechos, pero no establece ningún mecanismo de contralor.
\end{important}

\section{Trabajo forzoso (artículo 8)}

El artículo 8 de la Declaración reconoce el derecho a un trabajo libremente elegido y establece los compromisos de los Estados frente al trabajo forzoso.

\begin{quote}
«Toda persona tiene derecho a un trabajo libremente elegido y a ejercer cualquier oficio o profesión de acuerdo con las disposiciones nacionales vigentes. Los Estados Partes se comprometen a adoptar las medidas necesarias para eliminar toda forma de trabajo forzoso u obligatorio, exigido a un individuo bajo amenaza de sanción y para el cual no se haya ofrecido espontáneamente.»

Apartado 3: «Los Estados Partes se comprometen asimismo a adoptar medidas para garantizar la abolición de toda utilización de mano de obra que propicie, autorice o tolere el trabajo forzoso u obligatorio.»

Apartado 4: «Los Estados Partes se comprometen de modo especial a suprimir toda forma de trabajo forzoso u obligatorio o degradante que pueda ser utilizado: a) como medio de coerción o de educación política o como castigo por no tener o expresar el trabajador determinadas opiniones políticas o por manifestar oposiciones ideológicas al orden político, social, económico establecido; b) como método de movilización y utilización de la mano de obra con fines de fomento económico; c) como medida de disciplina en el trabajo; d) como castigo por haber participado el trabajador en actividades sindicales o huelgas; y e) como medida de discriminación racial, social, nacional, religiosa o de otra naturaleza.»
\end{quote}

La norma prohíbe el trabajo forzoso, pero la responsabilidad de adoptar las acciones necesarias se traslada a cada Estado: se lo toma como principio general, pero no hay una norma uniforme, sino \textbf{armonía legislativa} y dependencia de la decisión del Estado.

Más allá de la prohibición de la esclavitud, la definición alcanza prácticas que se observan en la sociedad argentina:

\begin{example}{trabajo forzoso en la práctica argentina (afiliación sindical)}
Para obtener determinado empleo hay que estar afiliado a cierto sindicato; en el caso de los camioneros, para poder trabajar deben estar afiliados al sindicato y tener determinadas tarjetas, pues de lo contrario no pueden circular. Se señala que ello obliga a afiliarse a un gremio para obtener trabajo, lo cual constituye trabajo forzoso.
\end{example}

Otras prácticas mencionadas son la participación obligada en manifestaciones políticas o movilizaciones para conservar el empleo (por ejemplo, cuando las intendencias movilizan). Se advierte una incoherencia: se afirma que no puede imponerse trabajo forzoso en las cárceles y, sin embargo, se permite obligar a asistir a manifestaciones políticas o a asociarse a determinados gremios para condicionar el acceso al trabajo.
% TODO: contenido incompleto o ambiguo en las notas originales (pasaje sobre trabajo en las cárceles).
Existen formas de trabajo forzoso que no se reducen al trabajo en negro, y se trata de un compromiso asumido por los Estados que los propios Estados no estarían cumpliendo.

\begin{important}
La Declaración distingue \textbf{derechos individuales} y \textbf{derechos colectivos}. Entre los colectivos se encuentran el sindicalismo y el derecho gremial (derecho de asociación); como derecho individual se reconoce el derecho a \textbf{asociarse o a no asociarse}.
\end{important}

\section{Trabajo infantil y protección del trabajo adolescente}

La Declaración regula dos cuestiones: evitar el trabajo infantil y proteger el trabajo adolescente. Como medida de protección de los menores se mencionó que no se permite extender la jornada laboral, ni siquiera mediante horas extra.

Respecto de qué se considera trabajo infantil, se observó que la Declaración \textbf{no fija una edad mínima}: solo indica que esta no debe ser inferior a aquella en que comienza la escolaridad obligatoria, y que se considera infantil el trabajo mientras dure la educación obligatoria.
% TODO: contenido incompleto o ambiguo en las notas originales (se menciona la edad de 18 años sin precisar a qué categoría se refiere; no se distingue con claridad entre infantil y adolescente).
La educación obligatoria depende de las normas nacionales. Por ello, en un país donde la obligatoriedad se limitara a la escuela primaria (en el ejemplo dado, 7 años, de modo que a los 13 años se dejaba de ser infante y se pasaba a ser adolescente) el límite sería menor. Como la Argentina extendió la obligatoriedad a la educación secundaria, para la Argentina el trabajo infantil se extendería hasta que se termina la secundaria (normalmente a los 17 o 18 años).

Las observaciones críticas formuladas son:
\begin{itemize}
    \item existe una \textbf{asimetría} entre lo que regula un Estado y lo que regula otro, y no se sabe si es favorable dentro de un ámbito de integración que haya esa diferencia; se señala que 3 o 5 años de diferencia no pesan igual a los 30 años que a los 10;
    \item la Declaración no define qué es un infante ni qué es un adolescente, que serían la base común para todos los Estados;
    \item condiciona el trabajo a que se termine de estudiar, con la finalidad de priorizar la educación frente al trabajo, pero desconoce una realidad: muchas veces los niños deben ayudar a sus padres en tareas rurales (ordeñar, recoger algodón) sin que ello les impida asistir a la escuela, y se trata de grupos numerosos en Latinoamérica;
    \item no hay una medida concreta al respecto.
\end{itemize}

\section{Garantía, promoción y respeto}

En la mayoría de las disposiciones generales y fundamentos, la Declaración no dice solamente que el Estado garantiza, sino que \textbf{asume el compromiso de promover}. Se trata, entonces, de obligaciones de garantía, de promoción y de respeto.

\begin{quote}
Artículo 3: «Los Estados Partes se comprometen a promover el desarrollo sostenible en la región, estimular la creación y el desarrollo de empresas sostenibles, promover el crecimiento de los mercados internos y regional y el fortalecimiento de la competitividad de las empresas sostenibles para el acceso a los mercados internacionales.»
\end{quote}

Se cita también un artículo (el 31, según la numeración indicada en clase) conforme al cual los Estados se comprometen a respetar los derechos contenidos en la Declaración y a promover su aplicación, así como la de las convenciones internacionales ratificadas.
% TODO: contenido incompleto o ambiguo en las notas originales (la docente indicó que ese artículo no era el más apropiado para ilustrar la promoción).

Se concluye que, aunque se trate de una declaración y no de una convención, hay un compromiso de los Estados de promover con actos concretos una economía sustentable que mejore la calidad de vida; la Declaración no dice cómo, pero señala el foco: «la zanahoria» hacia la cual se orienta el avance. Se menciona además que los Estados se comprometen a cumplir toda la normativa de la OIT, lo que se considera un compromiso muy escaso, pues se trata de convenciones ya firmadas y en un ámbito regional correspondería avanzar más. La Declaración es importante porque reconoce y promueve, más allá de garantizar.

Se señalan como contenidos de promoción: no discriminar entre hombres y mujeres, no discriminar por discapacidad y facilitar el acceso a la información; además, la capacitación (educación en el trabajo) crea obligaciones de promoción para los Estados y las empresas. Por ejemplo, los talleres gratuitos de oficios que brinda el Estado (electricidad, automotriz) cumplirían con el deber de capacitar para desarrollar una actividad remunerada, en relación de dependencia o en forma independiente.

\section{«Trabajo decente» y concepto de trabajador}

La Declaración \textbf{no define qué es trabajo} ni quién es trabajador, aunque se refiere al «trabajo decente». En la Unión Europea, en cambio, sí está definido el concepto de trabajador.

\begin{definition}{trabajador (Unión Europea)}
Es trabajador quien desarrolla una actividad bajo relación de dependencia, en forma continua y con un pago (en forma asalariada). Un profesional autónomo o el titular de una pyme, al no estar en relación de dependencia, no encuadra en el concepto.
\end{definition}

\begin{table}[H]
\centering
\small
\renewcommand{\arraystretch}{1.25}
\caption{«Trabajador» y «trabajo decente» en el Mercosur y en la Unión Europea}
\begin{tabularx}{\textwidth}{L{0.20\textwidth}YY}
\toprule
\textbf{Aspecto} & \textbf{Mercosur (Declaración Sociolaboral)} & \textbf{Unión Europea} \\
\midrule
\textbf{Concepto de trabajador} & No define qué es trabajo ni quién es trabajador. & Es trabajador quien desarrolla una actividad bajo relación de dependencia, en forma continua y asalariada. \\
\textbf{«Trabajo decente»} & Lo menciona, pero no define qué es «decente»; queda librado a la moral y las buenas costumbres de cada Estado. & Nunca habla de trabajo decente. \\
\textbf{Profesional autónomo o pyme} & No queda claro si está comprendido. & No está en relación de dependencia; no encuadra en el concepto de trabajador. \\
\textbf{Consecuencia ilustrada} & La Declaración se aplica solo al trabajo decente: no cubriría, por ejemplo, la actividad del traficante de drogas. & El concepto delimita con precisión a quién se aplican las normas. \\
\bottomrule
\end{tabularx}
\end{table}

Se ejemplifica con actividades que suponen una remuneración pero no serían trabajo decente: el ejercicio de la prostitución (se plantea que, así, la Declaración no regularía ni siquiera los casos de niñas de 12 años) y el tráfico de drogas. Se critica que la norma se refiera a la decencia sin definir qué es trabajo ni qué significa «decente», y que quede librado a la moral y las buenas costumbres de cada Estado.

Sí se refiere la Declaración al \textbf{trabajo sustentable}: promueve la actividad y la economía sustentables, aunque no dice cómo.

\section{La técnica: armonía legislativa y uniformidad legislativa}

La Declaración reenvía siempre a la normativa de los Estados Partes (cuando habla de no trabajo infantil, de igualdad de trato, de no discriminación, de acceso en igualdad de condiciones, la fórmula es «de conformidad con la normativa de los Estados Partes»). Por lo tanto, el sistema es de \textbf{armonía legislativa} y no de uniformidad legislativa.

\begin{definition}{armonía legislativa}
Remite, mediante una norma de conflicto, a la ley interna del Estado. El contenido lo define cada Estado. Ejemplo: se es «infante» mientras no se haya cumplido la educación obligatoria del propio Estado.
\end{definition}

\begin{definition}{uniformidad legislativa}
Establece una regla idéntica para todos los Estados del bloque. Ejemplo dado: «la edad mínima para el trabajo es de 18 años», igual para todos.
\end{definition}

\begin{table}[H]
\centering
\small
\renewcommand{\arraystretch}{1.25}
\caption{Armonía legislativa y uniformidad legislativa}
\begin{tabularx}{\textwidth}{L{0.20\textwidth}YY}
\toprule
\textbf{Aspecto} & \textbf{Armonía legislativa} & \textbf{Uniformidad legislativa} \\
\midrule
\textbf{Qué hace la norma} & Remite, mediante una norma de conflicto, a la ley interna de cada Estado. & Fija una regla idéntica para todos los Estados del bloque. \\
\textbf{Quién define el contenido} & Cada Estado, según su propia normativa. & La propia norma del espacio de integración. \\
\textbf{Ejemplo dado en clase} & «Infante» es quien no cumplió la educación obligatoria de su Estado; los Estados se comprometen a adoptar «las medidas necesarias» contra el trabajo forzoso. & «La edad mínima para el trabajo es de 18 años»: igual para todos. \\
\textbf{En la Declaración Sociolaboral} & Es la técnica que predomina: «de conformidad con la normativa de los Estados Partes». & No es la técnica adoptada. \\
\bottomrule
\end{tabularx}
\end{table}

\section{Órgano de seguimiento: la Comisión Sociolaboral}

La comisión sociolaboral actúa como organismo de contralor porque tiene la tarea de realizar el seguimiento y los informes (para conocer la situación de cada Estado) y de hacer propuestas de reformas o de avances legislativos y elevarlas al Grupo Mercado Común para que las sancione.

\begin{important}
El seguimiento es de tipo \textbf{informativo}: es un estado de situación. No es un seguimiento como el de la comisión de la OEA, que vigila e impone una sanción o un apercibimiento para que se revean las normas. La comisión puede poner en evidencia si en la Argentina o en Venezuela se cumple o no la Declaración, pero ningún organismo del Mercosur impone al Estado la exigencia de una actividad concreta para mejorar lo que no cumple.
\end{important}

\section{El valor agregado del trabajo (plusvalía) y la Constitución Nacional}

La Declaración habla de promover el desarrollo del \textbf{valor agregado del trabajo}.

\begin{definition}{valor agregado del trabajo}
Es la plusvalía (plusvalor) que el trabajo incorpora al producto.
\end{definition}

Desde una visión de izquierda, cuando una economía se vale del trabajo del trabajador para agregar valor a la producción lo está explotando, y ello solo beneficiaría al empresario. Sin embargo, la Declaración dice lo contrario: es importante el valor que el trabajo agrega, porque el país crece cuando lo que exporta no son solo commodities o materias primas, sino productos elaborados, que llevan valor agregado o plusvalía (trabajo de elaboración).

\begin{example}{exportar commodities frente a productos elaborados}
Una cosa es vender trigo y otra vender barritas de cereales: en el trigo hay trabajo (la siembra, la cosecha), pero no hay una elaboración que sume valor al producto final. Quien vende commodities necesita comprar producción final, más cara por su valor agregado, y la balanza económica le resulta desfavorable: vendió materia prima y compró producción elaborada.
\end{example}

\begin{example}{materia prima y manufactura}
Una cosa es vender lana (o alpaca, vicuña) y comprar después el pullover de cashmere inglés, y otra es manufacturar y vender el producto manufacturado, que implica incorporación de plusvalía (o valor agregado, según se lo mire).
\end{example}

Se concluye que, para crecer económicamente, hay que promover el trabajo sustentable, pero con valor incorporado: el progreso necesita el trabajo y el valor que este incorpora al producto, y la balanza de pagos se equilibra más con la venta de producto final que con la de commodities. Los gobernantes que firmaron la Declaración no eran, se señala, muy de derecha. Se promueve la industria y la comercialización industrializada con valor agregado, pero no se garantiza ni se obliga: la promoción tiene un objetivo, «una zanahoria delante». La Declaración combina garantía, promoción y respeto, porque quien promueve también respeta: promueve algo sin obligar.

\subsection*{Constitución Nacional, artículo 14: «industria lícita» frente a «industria útil»}

El proyecto y las Bases de Alberdi decían «trabajar y ejercer toda industria \textbf{útil}», mientras que la Constitución dice «trabajar y ejercer toda industria \textbf{lícita}». La Constitución fue más liberal: no importa si la actividad sirve o no, mientras sea legal. El texto de Alberdi, en cambio, estaba orientado al crecimiento de la nación y a un fin más social: había derecho a hacer todo lo que la nación requiriera para su crecimiento, en un momento en que el país se tenía que formar y crecer. La Constitución, por su parte, admitió todo lo que no fuera ilícito. Se señala que la diferencia muestra la mentalidad de cada uno: que vengan todos los que quieran hacer, mientras sea lícito, frente a promover lo que dé crecimiento.

\section{Derechos individuales y colectivos; el descanso dominical}

La Declaración enuncia derechos individuales (acceso al trabajo, igualdad de trato, igual remuneración por igual tarea) y derechos colectivos (derecho a la asociación o a la no asociación). Entre los derechos laborales figura el tope diario de horas, conforme a las convenciones internacionales, y el \textbf{descanso semanal}, con mención especial del \textbf{descanso dominical}.

Se plantea por qué no se dejó solo el descanso semanal y se agregó especialmente el domingo. La explicación expuesta es que la sociedad argentina se organizó sobre bases católicas y cristianas: el domingo era dedicado a la Iglesia y era el día de descanso para ir a misa, y esa tradición se conservó. La Declaración, de 2015, lo mantuvo no necesariamente con la finalidad de que se vaya a la iglesia, sino porque la sociedad está organizada para que el domingo se descanse (no hay clases y las oficinas públicas normalmente no trabajan).

Además, un día común favorece la reunión familiar: si cada miembro de la familia descansara un día distinto, el domingo no se encontraría nunca la familia reunida. Poner uniformidad en un día, especialmente el domingo en la medida de lo posible, tiende a conservar la unión familiar, basada en una tradición histórica vinculada con la religión.

\begin{note}
Se sostiene que, antes de criticar una costumbre o de cambiarla, conviene mirar cómo se fue evolucionando y a qué se apunta con ella: no solo a ir a misa, sino también a la reunión familiar. Se afirma que los países que descuidaron la familia son los que más conflictos tienen.
\end{note}

% ------------------------------------------------------------
% CORNELL CAPÍTULO 2
% ------------------------------------------------------------
\section{Cuadro Cornell del capítulo}

\begin{longtable}{L{0.27\textwidth}L{\dimexpr0.73\textwidth-4\tabcolsep\relax}}
\toprule
\textbf{Preguntas / palabras clave} & \textbf{Notas / respuestas} \\
\midrule
\endfirsthead
\toprule
\textbf{Preguntas / palabras clave} & \textbf{Notas / respuestas} \\
\midrule
\endhead

\textbf{¿Cuál es la norma laboral central del Mercosur?} &
La Declaración Sociolaboral del Mercosur (2015), elaborada por la comisión de asuntos sociolaborales del Grupo Mercado Común. Es una declaración: una «expresión de buenos deseos» sin mecanismo de contralor sancionatorio. \\[0.6em]

\textbf{¿Armonía o uniformidad legislativa? ¿Cuál adopta la Declaración?} &
Armonía: remite a la ley interna de cada Estado («de conformidad con la normativa de los Estados Partes»). Uniformidad sería una regla única (por ejemplo, edad mínima de 18 años). La Declaración funciona por armonía. \\[0.6em]

\textbf{¿Qué es el trabajo forzoso según el artículo 8 y qué crítica se hizo?} &
Se prohíbe como coerción política, método de movilización económica, medida de disciplina, castigo por huelga o discriminación. Se señalaron prácticas actuales contradictorias: afiliación sindical obligatoria para conseguir empleo y movilizaciones para conservarlo. La acción queda a cargo de cada Estado. \\[0.6em]

\textbf{¿Cómo trata la Declaración el trabajo infantil?} &
Protege contra el trabajo infantil y adolescente, pero no fija edad mínima: remite a la educación obligatoria de cada Estado. Genera asimetrías entre Estados y desconoce la realidad de familias que necesitan que los hijos colaboren. \\[0.6em]

\textbf{¿Qué diferencia hay entre obligaciones de garantía, promoción y respeto?} &
Respetar (artículo 31, según la numeración indicada) y promover (artículo 3: desarrollo sostenible, empresas sostenibles, mercados internos) son compromisos estatales que fijan un norte («la zanahoria») sin decir cómo lograrlo. \\[0.6em]

\textbf{¿Qué problema plantea el «trabajo decente» y cómo se compara con la UE?} &
El Mercosur lo usa sin definir «decente» ni «trabajador». La UE sí define trabajador (dependencia, continuidad, remuneración) y no habla de trabajo decente. \\[0.6em]

\textbf{¿Qué órgano controla el cumplimiento y con qué alcance?} &
La comisión sociolaboral: se reúne al menos dos veces por año, elabora informes y estadísticas y propone cambios por consenso. Su seguimiento es informativo; no impone sanciones (a diferencia de la comisión de la OEA). \\[0.6em]

\textbf{¿Qué es el valor agregado y por qué importa económicamente?} &
Es la plusvalía que el trabajo incorpora al producto. Exportar productos elaborados (no solo commodities) equilibra la balanza de pagos; la Declaración promueve el valor agregado. \\[0.6em]

\textbf{¿Qué diferencia hay entre «industria útil» y «industria lícita»?} &
Alberdi proponía industria útil, orientada al crecimiento de la nación; la Constitución, más liberal, admite toda industria lícita. \\[0.6em]

\textbf{¿Por qué la Declaración prevé el descanso dominical?} &
Por la base cristiana de la sociedad argentina y, sobre todo, porque un día común favorece la reunión familiar. Antes de cambiar una costumbre hay que preguntarse a qué apunta. \\

\midrule
\multicolumn{2}{L{\dimexpr\textwidth-2\tabcolsep\relax}}{\textbf{Resumen:} en el Mercosur, la Declaración Sociolaboral enuncia aspiraciones (no discriminar, erradicar el trabajo forzoso e infantil, promover el trabajo sustentable con valor agregado), pero lo hace mediante armonía legislativa, remitiendo a cada Estado y con un control meramente informativo.} \\
\bottomrule
\end{longtable}

% ============================================================
% CAPÍTULO 3
% ============================================================
\chapter[Derecho laboral: UE]{Derecho laboral en la Unión Europea}

\section{Libre circulación de trabajadores y entrada para buscar empleo}

La libre circulación es una fuente fundamental de la Unión Europea e incluye la libre circulación de trabajadores. Se plantea como pregunta capciosa de verdadero o falso la afirmación de que los tratados no hablaban de libertad de tránsito en un primer momento: siempre hablaron de ella, porque es un pilar básico.
% TODO: contenido incompleto o ambiguo en las notas originales (el pasaje oral sobre esta pregunta de verdadero o falso es confuso).

La regulación específica de la circulación laboral aparece a partir de 1957, con la formación de la Comunidad Económica Europea (Tratado de Roma). En 1951 se había formado la CECA; Euratom (Comunidad Europea de la Energía Atómica) no tiene que ver con la circulación de las personas. La Comunidad Económica Europea habla del mercado y del libre tránsito de personas, mercaderías, servicios, bienes y moneda única. La libre circulación y la igualdad de trato eran principios desde los orígenes y se sobreentendían; la regulación específica es de 1957.

Hay igualdad de trato de todos los ciudadanos europeos: se puede ingresar a un país para \textbf{buscar trabajo}, y no solamente con el trabajo ya conseguido. El régimen es el siguiente:
\begin{itemize}
    \item Se puede ingresar como turista o para buscar trabajo; si se declara que se viene a buscar trabajo, no se puede negar la entrada.
    \item Se puede buscar trabajo hasta \textbf{3 meses}.
    \item Después de los 3 meses hay que demostrar que se consiguió trabajo (tramitar la tarjeta de residencia laboral, para realizar una actividad económica bajo relación de dependencia) o demostrar medios suficientes, para no ser una carga para el Estado.
    \item No se puede pedir visa ni permiso de residencia; sí lo que tiene que ver propiamente con el trabajo.
\end{itemize}

La lógica es que la libre circulación y residencia se regulan de manera acorde con la libertad de trabajo.

\section{Igualdad de trato: armonía y no uniformidad}

El trabajador tiene el derecho a la igualdad de trato y al acceso al trabajo en igualdad de condiciones que un nacional del Estado de acogida. Aquí también rige la \textbf{armonía y no la uniformidad}: si al nacional se le exige un título habilitante para ocupar un trabajo, se lo pueden exigir al nacional de otro Estado.

Habría uniformidad si se exigiera exactamente lo mismo en todos los países (por ejemplo, tener 30 años o un título universitario para ingresar al Estado). Hay armonía porque se atiende a lo que cada Estado requiera: en uno puede pedirse título terciario para trabajar en una oficina pública y en otro título secundario. Lo que no se puede es distinguir entre un nacional y un extranjero dentro de la Unión Europea.

Tampoco se puede distinguir en cuanto al título: existe una tabla y un parámetro regulado en otro lado, con las reválidas y el reconocimiento de títulos (si el título es exactamente equivalente o si hay que hacer alguna reválida o control, lo que depende del título). Normalmente no se distingue si el título se obtuvo en el país de acogida o en el país de origen, y tiene la misma validez.

\section{Cupos y cargos reservados a nacionales}

No puede existir, en lo privado y tampoco en lo público, un cupo excluyente o para nacionales solamente (por ejemplo, que tenga que haber un 60\,\% o un 10\,\% de extranjeros). No puede existir un cupo en el empleo público, salvo para determinados empleos que pueden estar reservados a los nacionales, como los vinculados con la \textbf{seguridad nacional}. Mientras más alto sea el cargo, más puede ser que se excluya a los extranjeros. Esto no se considera discriminatorio en tanto sea lógico y no desproporcionado.

\begin{example}{el portero y el ingeniero de satélites (cargos reservados a nacionales)}
No puede fijarse un cupo máximo (por ejemplo, del 3\,\%) de extranjeros para quienes atienden el teléfono, pues pueden ser nacionales o extranjeros. En cambio, es válido que los cargos que organizan la seguridad nacional estén reservados a los nacionales: por ejemplo, los ingenieros de la oficina de satélites, si el satélite va a vigilar la seguridad nacional sobre el propio territorio.
\end{example}

\begin{note}
\textbf{Para reflexionar:} se plantea si en una selección deportiva puede exigirse ser nacional del Estado, es decir, si ello violaría la igualdad de trato en la Unión Europea. Se observa que en los clubes de fútbol (y se menciona la selección de Francia) hay más extranjeros que nacionales, y que exigir esa nacionalidad iría en contra de la regulación del derecho al trabajo de la Unión Europea. La cuestión se deja planteada, sin debate.
\end{note}

\section{Requisitos admisibles: título, idioma, experiencia y nacionalidad}

No se puede exigir ningún requisito extra por ser extranjero. Sí se pueden exigir determinados requisitos, \textbf{en igualdad de condiciones y proporcionadamente}: la posesión de títulos, hablar fluidamente el idioma nacional y acreditar experiencia.

\begin{example}{el idioma como requisito proporcional}
Quien va a trabajar a Alemania puede ser obligado a hablar alemán, en una medida proporcional al trabajo que haga. Lo que no se puede exigir es que ese aprendizaje del idioma se haya adquirido en el territorio de acogida: si la persona estudió alemán en España y lo habla con fluidez, sirve. Si no tiene título pero demuestra que habla el idioma con fluidez, no pueden excluirla.
\end{example}

\begin{example}{la experiencia no puede exigirse en el país de acogida}
Quien ingresa buscando trabajo no puede ser obligado a acreditar 5 años de experiencia en un empleo nacional: puede acreditar los 15 años de experiencia en el mismo trabajo en su país.
\end{example}

La \textbf{nacionalidad} sí puede exigirse para determinados trabajos (seguridad, policía, diplomacia y relaciones exteriores), donde se puede descartar al extranjero. Es proporcional, pero depende del cargo. Se concluye que la no discriminación y la igualdad de trato también rigen en la evaluación de esos requisitos y que el régimen de la Unión Europea es bastante claro.

% ------------------------------------------------------------
% CORNELL CAPÍTULO 3
% ------------------------------------------------------------
\section{Cuadro Cornell del capítulo}

\begin{longtable}{L{0.27\textwidth}L{\dimexpr0.73\textwidth-4\tabcolsep\relax}}
\toprule
\textbf{Preguntas / palabras clave} & \textbf{Notas / respuestas} \\
\midrule
\endfirsthead
\toprule
\textbf{Preguntas / palabras clave} & \textbf{Notas / respuestas} \\
\midrule
\endhead

\textbf{¿Cómo se regula el trabajo en la UE y desde cuándo?} &
La libre circulación y la igualdad de trato son principios desde los orígenes; la regulación específica comienza con la Comunidad Económica Europea (Tratado de Roma, 1957). \\[0.6em]

\textbf{¿Qué ocurre con quien ingresa a la UE para buscar trabajo?} &
Puede ingresar sin visa y buscar trabajo hasta 3 meses; luego debe tramitar la tarjeta de residencia laboral o demostrar medios suficientes, para no ser una carga para el Estado. \\[0.6em]

\textbf{¿Rige armonía o uniformidad en la igualdad de trato dentro de la UE?} &
Armonía: el extranjero comunitario tiene los mismos derechos que el nacional y se le pueden exigir los requisitos (por ejemplo, un título habilitante) que cada Estado exija a sus nacionales. Uniformidad sería exigir exactamente lo mismo en todos los países. \\[0.6em]

\textbf{¿Qué cargos pueden reservarse a los nacionales?} &
No puede haber cupos de extranjeros en el empleo privado ni en el público, salvo cargos como los vinculados con la seguridad nacional, la policía o la diplomacia, siempre que la exclusión sea lógica y no desproporcionada. \\[0.6em]

\textbf{¿Qué requisitos puede exigir un Estado de la UE a un trabajador de otro Estado?} &
Título, idioma (en forma proporcional) y experiencia, pero sin exigir que se hayan obtenido en el país de acogida y sin cupos de extranjeros. La nacionalidad solo en cargos como seguridad nacional, policía o diplomacia. \\

\midrule
\multicolumn{2}{L{\dimexpr\textwidth-2\tabcolsep\relax}}{\textbf{Resumen:} en la Unión Europea, la libre circulación y la igualdad de trato con los nacionales están reguladas desde 1957, con reglas claras sobre ingreso, residencia y requisitos admisibles (título, idioma, experiencia), sin cupos de extranjeros, salvo cargos reservados por razones de seguridad nacional.} \\
\bottomrule
\end{longtable}

% ============================================================
% CAPÍTULO 4
% ============================================================
\chapter[Repaso del parcial]{Consultas y repaso para el parcial}

\section{Acuerdo de prórroga y sumisión al foro}

\begin{definition}{sumisión al foro}
Siempre es \textbf{post litem}: implica que el actor interpone la demanda y que el demandado la acepta, de modo que ambos se someten a un juez que no es competente y le piden que intervenga y resuelva. No se interpone una demanda si no hubo un conflicto.
\end{definition}

\begin{definition}{acuerdo de prórroga}
Puede ser \textbf{pre litem} o \textbf{post litem}. Antes del conflicto se prevé que, si lo hubiera, las partes se someterán a determinado juez o árbitro; una vez ocurrido el conflicto, las partes acuerdan someter ese inconveniente a un juez o a un árbitro.
\end{definition}

\begin{example}{cláusula de prórroga en un contrato}
«Cualquier cuestión relativa a la interpretación, ejecución, derechos y obligaciones emergentes del presente contrato será resuelta por el tribunal de Mozambique, el tribunal de Uruguay o el árbitro Pepe.»
\end{example}

\begin{table}[H]
\centering
\small
\renewcommand{\arraystretch}{1.25}
\caption{Acuerdo de prórroga y sumisión al foro}
\begin{tabularx}{\textwidth}{L{0.20\textwidth}YY}
\toprule
\textbf{Aspecto} & \textbf{Acuerdo de prórroga} & \textbf{Sumisión al foro} \\
\midrule
\textbf{Momento} & Puede ser pre litem (antes del conflicto) o post litem (después). & Siempre post litem: no se demanda si no hay conflicto. \\
\textbf{Cómo se manifiesta} & Las partes eligen juez o árbitro, por cláusula contractual o por acuerdo posterior. & El actor interpone la demanda y el demandado la acepta: ambos se someten a un juez que no es competente. \\
\textbf{Ejemplo} & Cláusula: cualquier cuestión sobre interpretación, ejecución, derechos y obligaciones será resuelta por el tribunal elegido o por un árbitro. & Las dos partes le piden a ese juez que intervenga y resuelva. \\
\bottomrule
\end{tabularx}
\end{table}

\section{Temario del parcial: qué entra y qué no}

\begin{important}
\textbf{Insolvencia entra; sociedades no entra.} La insolvencia se evalúa mediante preguntas de verdadero o falso y de opción múltiple (no está para desarrollar). Sobre sociedades solo se habló de las sociedades multinacionales a propósito del caso Swift; como no se dio completo, se lo excluyó.
\end{important}

Las preguntas a desarrollar estarán orientadas a:
\begin{itemize}
    \item la \textbf{jurisdicción} en el Protocolo de Buenos Aires;
    \item la \textbf{ley aplicable} en el Reglamento de Roma de la Unión Europea (ley aplicable, no autonomía de la voluntad).
\end{itemize}

\section{Insolvencia internacional: Tratados de Montevideo y Unión Europea}

\subsection{Sistemas y juez competente}

En el Mercosur no hay una norma uniforme sobre insolvencia: se aplican los Tratados de Montevideo; Bolivia y Venezuela se rigen por el Código Sánchez de Bustamante. Hay normativas dispersas, pero todas receptan la extraterritorialidad.
% TODO: contenido incompleto o ambiguo en las notas originales (la docente duda entre extraterritorialidad y territorialidad al referirse a estas normativas).
Los Tratados de Montevideo hablan de universalidad o extraterritorialidad, ya sea con procedimiento único o con la posibilidad de procedimientos plurales: se permite un procedimiento local siempre que se demuestre una sede totalmente independiente. La Unión Europea también regula el sistema mediante la extraterritorialidad y lo llama \textbf{procedimiento principal} y \textbf{procedimiento secundario}.

El juez competente para abrir el procedimiento principal es:
\begin{itemize}
    \item en los \textbf{Tratados de Montevideo}: el del \textbf{domicilio} del deudor;
    \item en la \textbf{Unión Europea}: el del \textbf{centro de intereses principales} del deudor.
\end{itemize}

\begin{important}
No debe confundirse el centro de intereses principales con el establecimiento: el establecimiento define el procedimiento secundario. El centro de intereses principales no está definido: es una cuestión más fáctica, y el juez evalúa si en ese lugar se encuentra el centro de actividades principales. Es una pregunta frecuente de verdadero o falso y de opción múltiple.
\end{important}

\begin{definition}{establecimiento (Unión Europea)}
Cualquier lugar donde se desarrolle una actividad no transitoria con recursos humanos y bienes, reciba o no órdenes de otra sede. Allí se abre el procedimiento secundario.
\end{definition}

\subsection{Procedimiento local o secundario}

\textbf{Tratados de Montevideo.} Si se detecta en otro Estado parte un establecimiento totalmente independiente (ni en las directrices ni en lo económico dependiente de la casa principal), los acreedores locales pueden pedir la apertura de un procedimiento local. Ese procedimiento se limita territorialmente a los acreedores, al deudor, a las deudas y a los bienes situados en ese territorio. Se requiere que el establecimiento sea autónomo en la toma de decisiones y en lo económico (con patrimonio propio) y que haya acreedores locales. Vencido el plazo de la publicación de edictos, hay 60 días para abrir el procedimiento local; no hace falta probar la cesación de pagos (el deudor puede estar al día con sus pagos). Esto no significa que se base en el principio de territorialidad: si así fuera, habría que demostrar el estado de insolvencia. Si no se abre el procedimiento local (porque no hay acreedores locales o estos no lo pidieron), todo el patrimonio tramita en el procedimiento del domicilio del deudor.

\textbf{Unión Europea.} El procedimiento secundario solo puede ser de liquidación de bienes y se rige por su propia ley, como en los Tratados de Montevideo. Se abre para los acreedores locales con los bienes locales, donde haya un establecimiento, y no distingue si este es dependiente o independiente.

\subsection{Cuadro comparativo}

\begin{longtable}{L{0.21\textwidth}L{\dimexpr(0.79\textwidth-6\tabcolsep)/2\relax}L{\dimexpr(0.79\textwidth-6\tabcolsep)/2\relax}}
\caption{Insolvencia internacional: Tratados de Montevideo y Unión Europea}\\
\toprule
\textbf{Aspecto} & \textbf{Tratados de Montevideo} & \textbf{Unión Europea} \\
\midrule
\endfirsthead
\toprule
\textbf{Aspecto} & \textbf{Tratados de Montevideo} & \textbf{Unión Europea} \\
\midrule
\endhead

\textbf{Juez del procedimiento principal} & El del domicilio del deudor. & El del centro de intereses principales del deudor (cuestión más fáctica: el juez evalúa dónde está; no se lo define). \\[0.4em]
\textbf{Procedimiento local o secundario} & Puede abrirse otro procedimiento autónomo donde haya un establecimiento totalmente independiente de la casa principal (en decisiones y en lo económico) y acreedores locales. & Se abre donde exista un establecimiento; solo puede ser de liquidación de bienes. No distingue entre dependiente e independiente. \\[0.4em]
\textbf{Concepto de establecimiento} & Independiente de la casa principal en la toma de decisiones y con patrimonio propio. & Cualquier lugar donde se desarrolle una actividad no transitoria con recursos humanos y bienes, reciba o no órdenes. \\[0.4em]
\textbf{Alcance del procedimiento local} & Acreedores, deudas y bienes situados en ese territorio; no hace falta probar la cesación de pagos. Tras el vencimiento de los edictos hay 60 días para abrirlo. & Acreedores locales con bienes locales. \\[0.4em]
\textbf{Si no se abre el procedimiento local} & Todo el patrimonio tramita en el procedimiento del domicilio del deudor. & --- \\[0.4em]
\textbf{Base común} & Extraterritorialidad: juicio único o pluralidad de procedimientos. & Extraterritorialidad: procedimiento principal y secundario. \\
\bottomrule
\end{longtable}

\subsection{Aclaraciones}

\begin{itemize}
    \item \textbf{Tratado de Montevideo de 1940.} Aunque no se hubiera abierto el procedimiento secundario y todos los bienes fueran al procedimiento principal (el del domicilio del deudor), los acreedores locales conservaban la preferencia sobre los bienes situados en el territorio: se hablaba de juicio único con pluralidad de masas procesales, tramitado ante aquel juez y con aquella ley. En lo demás no había diferencias relevantes entre los dos Tratados de Montevideo, y no se pregunta en ese sentido.
    \item \textbf{Base común.} En lo fundamental, ambos sistemas se basan en la extraterritorialidad y permiten tanto un juicio único como la pluralidad de procesos. Se agrega que, para que haya juicio único, el establecimiento tiene que ser dependiente o haberse realizado actos aislados, y que, si tiene establecimiento dependiente, puede abrirse la posibilidad de pluralidad de procesos abriendo un procedimiento autónomo donde tiene ese establecimiento totalmente independiente.
    % TODO: contenido incompleto o ambiguo en las notas originales (el pasaje parece contradictorio respecto de establecimiento dependiente e independiente).
    \item \textbf{Prestación más característica.} No tiene nada que ver con la insolvencia ni con la jurisdicción: es una teoría para la ley aplicable. En jurisdicción, el lugar de cumplimiento se interpreta como cualquier lugar de cumplimiento de cualquiera de las obligaciones del contrato.
    \item \textbf{Antecedentes.} Se explicó la evolución inicial para mostrar si la quiebra se regulaba como sanción, con una idea de salvar al deudor o con una orientación al mercado, y así evidenciar las asimetrías entre legislaciones y las dificultades para acordar una ley uniforme en materia de insolvencia. El fundamento no se preguntará en el examen.
\end{itemize}

\section{Ley aplicable en el Reglamento de Roma}

A falta de autonomía de la voluntad, la ley aplicable en Roma se determina del siguiente modo. Un catálogo (artículo 3, según lo expuesto) indica, para ciertos contratos, la ley de la \textbf{residencia de la parte que tiene a su cargo el cumplimiento de la prestación más característica}, con excepciones (por ejemplo, los inmuebles). Si el contrato no está en esos casos, se aplica la ley del país donde se presenten los \textbf{vínculos más estrechos}, que es un criterio más flexible y no tan rígido; después se aclara que los vínculos más estrechos se presumen donde se encuentra la residencia de la parte que cumple la prestación más característica, con lo cual en cierta medida se vuelve a lo mismo.

\begin{important}
Las preguntas se orientan al criterio general de ley aplicable en materia de contratos en la Unión Europea: si es un criterio rígido o flexible, si hay un catálogo o no y cómo se resuelven los contratos que quedan fuera del catálogo. No se preguntó nada sobre el contrato de transporte.
\end{important}

\begin{table}[H]
\centering
\small
\renewcommand{\arraystretch}{1.25}
\caption{Ley aplicable a falta de autonomía de la voluntad}
\begin{tabularx}{\textwidth}{L{0.20\textwidth}YY}
\toprule
\textbf{Aspecto} & \textbf{Criterio rígido} & \textbf{Criterio flexible} \\
\midrule
\textbf{Quién determina la ley} & La propia ley (lugar de cumplimiento o de celebración). & El juez, según la ley que presente los vínculos más estrechos. \\
\textbf{Fundamento} & Voluntad presunta de las partes. & Análisis del caso concreto. \\
\textbf{Reglamento de Roma (según clase)} & Un catálogo de contratos (con excepciones, por ejemplo inmuebles) remite a la residencia de la parte que cumple la prestación más característica. & Si el contrato no encaja en el catálogo: ley del país con los vínculos más estrechos. \\
\bottomrule
\end{tabularx}
\end{table}

\section{Repaso general de los contenidos del parcial}

Se recuerdan los siguientes ejes del programa:
\begin{itemize}
    \item \textbf{Derechos humanos}: la Declaración Universal, el pacto económico y social; si se trata de normativa propia del Mercosur o de normativa internacional aplicable dentro del área, por haber sido incorporada como derecho autónomo o como derecho internacional.
    \item \textbf{Libertad de tránsito y residencia} en el Mercosur y en la Unión Europea, con los derechos y requisitos de cada plazo.
    \item \textbf{Derecho laboral} (visto en esta clase), con un par de fallos de la Unión Europea sobre derechos sociales (jubilación, asistencia, salud), tema del que se indica que debió hablarse un poco más.
    \item \textbf{Insolvencia}: sistema de la territorialidad, sistema de la extraterritorialidad, Tratados de Montevideo y reglamento de la Unión Europea.
    \item \textbf{Contratos}: formación o celebración (lugar y momento): teoría de la aceptación, de la remisión, de la recepción, del conocimiento y del conocimiento presunto; forma del contrato (validez extrínseca: si debe ser escrito, firmado o puede ser electrónico), que normalmente se rige por la ley del lugar de celebración; ley aplicable al contrato en sí (existencia, derechos y obligaciones, capacidad, licitud del objeto).
    \item \textbf{Autonomía de la voluntad}: principio absoluto o relativo; fundamento (si es inherente a la naturaleza humana, principio natural, o si depende del legislador); límites (orden público, internacionalidad objetiva); y, a falta de autonomía, ley aplicable con criterios rígidos o flexibles (se indica que hay pregunta a desarrollar), con la cuestión de cuál recepta Roma y cuál recepta el derecho argentino.
    \item \textbf{Jurisdicción}: directa, indirecta y prórroga; en el Mercosur, el Protocolo de Las Leñas y el Protocolo de Buenos Aires; en Europa, el Reglamento de Bruselas.
    \item \textbf{No entra}: sociedades.
\end{itemize}

\subsection{Libertad de tránsito y residencia en la Unión Europea}

\begin{definition}{grupo familiar (Unión Europea)}
Definición amplia, que comprende al cónyuge, la pareja de hecho, los hijos de ambos, los hijos de la pareja de hecho, los progenitores del ciudadano y los progenitores de la pareja de hecho. El derecho a la libre residencia lo tiene el ciudadano europeo y comprende también a su grupo familiar.
\end{definition}

\begin{table}[H]
\centering
\small
\renewcommand{\arraystretch}{1.25}
\caption{Libertad de tránsito y residencia en la Unión Europea}
\begin{tabularx}{\textwidth}{L{0.20\textwidth}L{0.19\textwidth}Y}
\toprule
\textbf{Estadía} & \textbf{Plazo} & \textbf{Qué se exige / qué derechos hay} \\
\midrule
\textbf{Estadía corta} & Hasta 3 meses & Se puede entrar como turista o para buscar trabajo; no pueden exigirse visa ni permiso de residencia. Quien dice que viene a buscar trabajo no puede ser rechazado por eso. \\
\textbf{Residencia de más de 3 meses} & Más de 3 meses & Hay que demostrar que se consiguió trabajo (tramitar la tarjeta de residencia laboral) o contar con medios suficientes para no ser una carga para el Estado. \\
\textbf{Residencia permanente} & Más de 5 años en el mismo Estado & Se adquiere el derecho de residencia permanente. El derecho de residencia alcanza también al grupo familiar. \\
\bottomrule
\end{tabularx}
\end{table}

Para cada una de estas libertades de tránsito y residencia conviene recordar qué derechos se tienen, qué papeles se piden y qué derechos se adquieren al ir superando los distintos plazos.

% ------------------------------------------------------------
% CORNELL CAPÍTULO 4
% ------------------------------------------------------------
\section{Cuadro Cornell del capítulo}

\begin{longtable}{L{0.27\textwidth}L{\dimexpr0.73\textwidth-4\tabcolsep\relax}}
\toprule
\textbf{Preguntas / palabras clave} & \textbf{Notas / respuestas} \\
\midrule
\endfirsthead
\toprule
\textbf{Preguntas / palabras clave} & \textbf{Notas / respuestas} \\
\midrule
\endhead

\textbf{¿Qué diferencia hay entre acuerdo de prórroga y sumisión al foro?} &
El acuerdo puede ser pre o post litem; la sumisión al foro siempre es post litem (el actor demanda y el demandado acepta al juez no competente). \\[0.6em]

\textbf{¿Qué entra y qué no entra en el parcial?} &
Insolvencia entra (verdadero o falso y opción múltiple); sociedades no entra. Las preguntas a desarrollar se orientan a la jurisdicción en el Protocolo de Buenos Aires y a la ley aplicable en el Reglamento de Roma. \\[0.6em]

\textbf{¿Cómo se determina el juez en la insolvencia internacional?} &
Tratados de Montevideo: domicilio del deudor, y procedimiento local solo con establecimiento independiente. UE: centro de intereses principales y procedimiento secundario (liquidación) donde haya establecimiento. \\[0.6em]

\textbf{¿Qué ley rige el contrato sin autonomía de la voluntad en Roma?} &
Un catálogo que remite a la residencia de la parte que cumple la prestación más característica y, si el contrato no encaja, la ley con los vínculos más estrechos. La prestación característica sirve para la ley aplicable, no para la jurisdicción. \\[0.6em]

\textbf{¿Qué plazos y derechos de residencia hay en la UE?} &
Hasta 3 meses (sin visa ni permiso de residencia); más de 3 meses (trabajo conseguido, con tarjeta de residencia laboral, o medios suficientes); más de 5 años en el mismo Estado, residencia permanente. El derecho alcanza al grupo familiar. \\

\midrule
\multicolumn{2}{L{\dimexpr\textwidth-2\tabcolsep\relax}}{\textbf{Resumen:} la clase cierra con el repaso del parcial: prórroga y sumisión al foro, insolvencia (Tratados de Montevideo frente a centro de intereses principales), ley aplicable en el Reglamento de Roma y los demás ejes del programa.} \\
\bottomrule
\end{longtable}

\end{document}
